\subsection{نتایج تجربی}
عملکرد تجربی ماژول \lr{dynamic\_rank} روی مدل‌های پایه بر مبنای ارزیابی‌های واقعی ثبت‌شده در نتایج آزمایشگاهی در جدول \ref{tab:benchmark_dynamic_rank} آورده شده است:

\begin{table}[H]
\centering
\small
\caption{نتایج بنچ‌مارک تجربی فشرده‌ساز \lr{dynamic\_rank} بر روی مدل‌های پایه}
\label{tab:benchmark_dynamic_rank}
\begin{tabular}{lccccc}
\toprule
\textbf{مدل پایه \lr{(Base Model)}} & \textbf{دقت \lr{(Accuracy)}} & \textbf{شاخص \lr{F1}} & \textbf{حافظه \lr{(Peak VRAM)}} & \textbf{توان \lr{(Throughput)}} & \textbf{وضعیت} \\
\midrule
\lr{DistilBERT} & \lr{90.60\%} & \lr{0.9099} & \lr{357.28 MB} & \lr{88,554.92 tok/s} & \lr{PASS} \\
\lr{RoBERTa} & \lr{94.15\%} & \lr{0.9425} & \lr{503.28 MB} & \lr{14,957.99 tok/s} & \lr{PASS} \\
\lr{GPT-2} & \lr{99.72\%} & \lr{0.0000} & \lr{552.93 MB} & \lr{12,097.77 tok/s} & \lr{PASS} \\
\lr{Qwen} & -- & -- & -- & -- & \lr{SKIP} \\
\bottomrule
\end{tabular}
\end{table}

این نتایج نشان می‌دهند که الگوریتم تخصیص پویا توانسته در مدل‌های انکودر (\lr{DistilBERT} و \lr{RoBERTa}) دقت بسیار بالای بالای ۹۰ درصد را همراه با گذر موفق از دروازه کیفیت (\lr{PASS}) حفظ نماید. در ارزیابی مدل \lr{GPT-2} آزمون مبتنی بر تکمیل متن بوده و دقت توکنی \lr{99.72\%} حاصل شده است. برای مدل \lr{Qwen} بنچ‌مارک به دلیل محدودیت‌های حافظه در سرور اجرا نشده و در وضعیت \lr{SKIP} قرار گرفته است.
